\begin{lemma}[Chordless cycles over a complete named base]
\label{lem:typed-complete-base-chords}
Let $D$ be a finite complete atomic typed diagram and
let $X$ be another disjoint finite set of formal vertices.
Suppose every $x\in X$ has a specified atomic edge to
EVERY $d\in D$, while the specified edges within $X$
are those of a graph $H$. Then each specified
\emph{chordless simple cycle} is either entirely contained
in $H$, or is a triangle meeting $D$.
In particular, to test individual old-cycle completability
it suffices to test all displayed triangles meeting $D$
and the chordless cycles of $H$.
\end{lemma}
\begin{proof}
A simple cycle meeting exactly one vertex $d$ of $D$
has all of its other vertices in $X$. If it has more
than three vertices, one of the internal $X$ vertices
is not a neighbour of $d$ ALONG the cycle, but is joined
to $d$ by a specified spoke: a chord.
If it meets two or more vertices of $D$, every two
of these vertices are adjacent by a specified old
base edge. Hence a chordless cycle can have at most
two base vertices; if two occur, they must be
consecutive along the cycle. Unless the cycle is a
triangle, the other path between them has an internal
vertex $x\in X$ not adjacent to at least one of those
base vertices ALONG the cycle. Its specified spoke
to that base vertex is a chord. Thus a nontriangular
chordless cycle cannot meet $D$.

For the final assertion, split a specified cycle along
any specified chord. Both resulting strictly shorter
cycles may be homomorphically realised with their
common chord given an endpoint colour; when the chord
is already specified, one can instead glue realizations
of the two resulting cycles along its complete coloured
two-vertex substructure using old SAP. Repeating the
splitting shows completability by induction on length.
\end{proof}

\begin{lemma}[Typed exact-coloured orbit infinitude]
\label{lem:typed-exact-orbit-infinitude}
Let \(p(x/D)\) be an exact atomic one-point type, with \(x\) in sort \(a\),
and, writing \(b_d\) for the unary sort of \(d\), let
\(t_d:a\to b_d\) be its spoke type to \(d\in D\).  Let \(J\) be a proper
represented meet-irreducible on \(a\).  If the \(J\)-class of \(x\) does
not belong to the full quotient-transition closure of \(D\), then
realisations of \(p\) occupy infinitely many \(J\)-classes.
\end{lemma}

\begin{proof}
The hypothesis forces the distinguished point $x$
to be outside the real base $D$: otherwise its
$J$-class would already be the named $J$-projection
of a member of $D$. Thus all old spokes $t_d$
below are ACTUAL DIVERSITY atoms, not logical equality.

For every \(d\in D\), put
\(P_d=\{t_d\}\starprod\{\bar t_d\}\), the
STRONG off-diagonal typed diversity profile of
the two-edge fork through the sort of \(d\). We have
\[
 P_d\not\subseteq J,
\tag{QT4}\label{eq:typed-unforced-spoke}
\]
Indeed otherwise the ACTUAL equality-inclusive RAW
product $t_d;t_d^{\smile}$ is contained in $J$:
its diversity restriction is exactly $P_d$
by original visible SAP injectivisation
\cref{lem:typed-injectivisation},
and its only remaining possible old atom is
the TRUE real diagonal, already in $J$.
Apply \cref{lem:typed-quotient-transport}
to the original OLD typed atomic relation
$t_d:a\to b_d$ and the proper represented
MI kernel $J$ on $a$. Its transported
kernel $J^{t_d}=t_d^{\smile};J;t_d$
on the sort of $d$ is again a represented
proper MI, and is NOT real equality, by
\cref{lem:no-real-algebraic-transition}.
The bijection of ACTUAL quotient sorts
$\pi_{t_d,J}:a/J\to b_d/J^{t_d}$
is one of the named terminal unary transitions,
and its inverse is named as well (the
transition for the converse atomic relation).
For the ACTUAL spoke $x\,t_d\,d$ this gives
the literal term identity
\[
 F_J(x)
   =\pi_{t_d,J}^{-1}
                \bigl(F_{J^{t_d}}(d)\bigr).
\]
The right-hand side is a NAMED unary
quotient-transition term value of the
old real point $d\in D$, contradicting
the hypothesis that $F_J(x)$ is outside
the FULL unary quotient-transition closure of $D$.
This proves (QT4) without silently
equating a diversity guard with a RAW
relation or assuming a target equality kernel.  Put
\[
 C_p=\bigcap_{d\in D}P_d,\qquad
 C_p=L_{aa}\text{ when }D=\varnothing.
\]
We must NOT claim that a STRONG diversity profile contains
the logical diagonal: it does not. Instead, realise the
COMPLETE old finite one-point type \(p(x/D)\) and
strongly amalgamate two copies over their full common
base \(D\). The two new points \(x,x'\) are distinct
and have equal prescribed spoke colours \(t_d\) to
every \(d\in D\). Their actual diversity colour
\(c=\col(x,x')\) therefore belongs to EVERY \(P_d\).
Thus \(C_p\ne\varnothing\) and indeed every
partial intersection is nonempty. In the empty-base
case \(C_p=L_{aa}\) is nonempty by SAP and
the existence of at least one same-sort diversity
atom.

If \(C_p\subseteq J\) and \(D\ne\varnothing\),
repeated use of \cref{lem:typed-prime-intersection}
on these nonempty STRONG exact domains \(P_d\)
would force some factor \(P_d\subseteq J\), contrary
to \eqref{eq:typed-unforced-spoke}.
For \(D=\varnothing\), the properness of \(J\)
already gives \(C_p\not\subseteq J\).
Choose an actual same-sort diversity atom
\(e\in C_p\setminus J\).

Fix any positive integer \(n\). Choose a finite graph
\(H\) with chromatic number \(>n\) and girth
\[
       g(H)>\max\{3,N_e+1\},
\]
where \(N_e\) is the TYPED STRONG same-sort
pure-power threshold from
\cref{lem:typed-terminal-pure-power}.
Take a disjoint formal vertex \(x_v\) of the fixed
sort \(a\) for each \(v\in V(H)\).
Keep the COMPLETE actual atomic diagram of \(D\);
give every \(x_v\) the SAME complete spoke type
\(p(x/D)\), and prescribe the old diversity
atom \(e\) exactly on graph edges \(x_vx_w\)
when \(vw\in E(H)\). All other clone-clone
edges are left unspecified.

We verify EVERY chordless specified simple cycle,
rather than merely asserting that long cycles are
harmless. By
\cref{lem:typed-complete-base-chords}, such a cycle
is either a triangle meeting the old base \(D\)
or is entirely within the clone graph \(H\).
A triangle with two base vertices is legal
because \(p(x/D)\) is a REALISED complete one-point
type; a triangle \((d,x_v,x_w)\) with one base
vertex is legal because
\(e\in \{t_d\}\starprod\{\bar t_d\}=P_d\).
A clone-only cycle has length
\(\ell>N_e+1\), so delete one edge as a
closing side; the remaining \(e\)-path
has length \(\ell-1\ge N_e\), and its
STRONG endpoint profile is exactly the
connectivity diversity palette \(F_e\).
Since \(e\in F_e\), the closing edge is
compatible. This gives an old realisation
of each clone-only cycle using SAP
injectivisation if the represented walk
initially repeats vertices. Every specified
simple cycle with chords is consequently
completable by the preceding chord lemma.

Semantic homomorphic CCO now yields ONE
homomorphism $f$ of the ENTIRE finite typed
partial graph into the old homogeneous source.
The complete OLD atomic diagram on $D$ is
specified, and every displayed pair atom
between two DISTINCT old base vertices is a
diversity atom. Thus $f|_D$ is an induced
embedding of the COMPLETE original base.
By ORIGINAL finite-diagram homogeneity,
choose ONE automorphism $h\in\Aut(\M)$
sending $f(d)$ to the ACTUAL pinned old
point $d$ for every $d\in D$; replace
$f$ by $h\circ f$. This common alignment
must precede orbit counting: otherwise
the clone images would realise the
one-point type over $f(D)$, not
literally over the stated base $D$.
For the orbit counting we need not force the
distinct formal clones to have distinct images;
however every specified graph edge is a
diversity atom, so adjacent vertices have
different images and, since \(e\notin J\),
different \(J\)-classes. Every clone image
realises the prescribed complete type
\(p(x/D)\) over the same base image.
Thus the \(J\)-classes of its realisations
properly colour \(H\), so at least
\(\chi(H)>n\) such classes occur.
In the support-normal finite-arity setting,
one additionally freezes the complete factors
\(D\cup\{x_v\}\) and uses the stated
support-normal completion interface;
no such interface is needed for this
binary theorem.
\end{proof}

\begin{corollary}[No hidden algebraicity outside unary closure of an original real base]
\label{cor:typed-no-real-base-hidden-algebraicity}
In the terminal visible SAP+semantic-CCO named-sorted source,
for every finite COMPLETE original real $D$,
\[
 \left\{u\in|\widehat\M|:
       |\Aut(\M/D)\cdot u|<\infty\right\}
                =D^{\coord}.
 \tag{PG-ACL}\label{eq:typed-no-real-base-hidden-algebraicity}
\]
Here $\widehat\M$ is the ACTUAL full interpreted hidden
structure, including the original real and all proper
represented quotient sorts, and $D^{\coord}$ is its full
point-generated unary closure of $D$.  Consequently
the orbit-theoretic algebraic closure of the original
real base $D$ in that full hidden structure has size
at most $m_{\mathrm{typ}}|D|$, with the empty-base
case interpreted literally. This asserts nothing
about algebraicity over an arbitrary orphan-only base.
\end{corollary}
\begin{proof}
Every vertex of $D^{\coord}$ is a value of a NAMED
unary term on a point of $D$, so is fixed by
$\Aut(\M/D)$.  Conversely take any hidden
$u\notin D^{\coord}$. By point generation of the
FULL interpreted hidden universe, write
$u=\tau(q)$ for an old real $q$ and one named
real-source unary term $\tau$.
The normal form \cref{lem:typed-coordinate-normal-form}
says either $\tau$ is the real identity, or it
has a represented proper meet-irreducible
original kernel $J$.
In the real case $q\notin D$ and visible SAP
gives arbitrarily many realisations of the
complete old one-point type $q/D$, hence
an infinite orbit under $\Aut(\M/D)$.
In the quotient case, if $[q]_J$ belonged
to the FULL unary quotient-transition closure of
$D$, then the normal form of $\tau$ would
put $u$ in $D^{\coord}$, a contradiction.
Thus \cref{lem:typed-exact-orbit-infinitude}
gives infinitely many $J$-classes realised by
the same old type $q/D$.  Old homogeneity
moves $q$ to each representative while
fixing $D$. The kernel biconditional of
$\tau$ makes the corresponding values of
$u$ pairwise distinct.  Therefore its
orbit is infinite. The uniform cardinality
bound is \cref{lem:typed-unary-closure-bound},
since the unary closure of finite real $D$
is the union of its one-generated closures.
\end{proof}

We apply the general point-generated SAP criterion of
\cref{lem:typed-point-generated-sap-criterion} to the preceding
coordinate-orbit calculation.


\begin{proposition}[Typed point-generated coordinate SAP]
\label{prop:typed-coordinate-sap}
Full point-generated coordinate expansions in the transition-enriched hidden
language have strong amalgamation over full point-generated common
substructures.
\end{proposition}

\begin{proof}
The source-specific orbit theorem
\cref{cor:typed-no-real-base-hidden-algebraicity}
gives the following condition for EVERY finite original
real base $D$ and EVERY actual point-generated hidden
coordinate vertex $u\notin D^{\coord}$:
\[
             |\Aut(\M/D)\cdot u|=\infty .
 \tag{TOS}\label{eq:typed-coordinate-orbit-stabilizer}
\]
This is exactly the hypothesis of the standalone
point-generated SAP criterion
\cref{lem:typed-point-generated-sap-criterion}.
That criterion applies because all named coordinate maps and
transition functions are unary and equivariant, every finite
old real factor has finite FULL hidden closure, and the
complete old real copy of a common base aligns by
old homogeneity.  Therefore the two FULL point-generated
hidden factors have one STRONG point-generated amalgam.

The hard ORIGINAL typed SAP+CCO content is precisely
the orbit-infinitude argument inside
\cref{lem:typed-exact-orbit-infinitude}; the generic
finite-coset step does NOT assert SAP over an orphan-only
hidden base. That case is proved separately below.
\end{proof}

\begin{lemma}[Meet-prime MI coordinate dichotomy inside one exact Helly fibre]
\label{lem:typed-mi-fibre-infinite-dichotomy}
Fix ONE original named transitive real sort $a$ of
the terminal symmetric-on-sorts SAP+homomorphic-CCO
source, with full invariant-lattice NR2.
Let $E_1,\ldots,E_m>\Id_a$ be original
invariant equivalences and let
$C_i$ be actual $E_i$-classes which
meet PAIRWISE (allow $m=0$).
Put
\[
       H=\bigcap_{i=1}^m E_i,\qquad
       X=\bigcap_{i=1}^m C_i,
\]
with $H=\Top_a$, $X=M_a$ for $m=0$.
Then $X$ is ONE INFINITE actual $H$-class.
For EVERY proper meet-irreducible original
invariant equivalence $J$ on the SAME
original sort there is the dichotomy
\begin{enumerate}[label=(\roman*)]
\item If $H\le J$, then some $E_i\le J$,
  and the actual $J$-class of any $x\in X$
  is already determined by that SINGLE
  original $E_i$-class. This case cannot
  arise when $m=0$ and $J<\Top$.
\item If $H\nleq J$, then $X$ meets
  INFINITELY MANY actual $J$-classes:
  $[H:H\cap J]=\infty$.
\end{enumerate}
Furthermore, given finitely many such
unforced MI kernels $J_1,\ldots,J_s$
and for each a FINITE list of forbidden
actual $J_k$-classes, ONE $x\in X$
can be chosen to avoid all lists
SIMULTANEOUSLY.
This statement concerns genuine
original class incidence, NOT a complete
fine old one-point spoke type, and does
not replace \cref{lem:typed-exact-orbit-infinitude}.
\end{lemma}
\begin{proof}
By \cref{lem:typed-terminal-all-equivalences-exact},
the FULL finite original invariant-equivalence
lattice on $a$ is distributive and every
invariant equivalence is FULL RAW exact.
For $m\ge1$, its original exact
equivalence classes $C_i$ are
pairwise-intersecting actual FULL RAW
exact balls. By
\cref{thm:full-typed-exact-balls}(iii)
they have a common original real
$x\in X$, whence
$X=[x]_H$. The latter is INFINITE
by original SAP and
\cref{lem:rh-sap-invariant-monolith},
since the finite meet $H$ of
nonidentity invariant equivalences
remains nonidentity. For $m=0$,
$X=M_a$ is infinite.

The exact meet-prime inference is
the already established finite-lattice
cofinality lemma
\cref{lem:foundational-MI-coordinate-cofinal}:
from $H=\bigcap_iE_i\le J$
for proper meet-irreducible $J$
one obtains $E_i\le J$ for
at least one $i$.
This is a REUSE of the earlier
abstract CRT geometry, not
a new lattice assertion.
The $J$-class of $x$ is then
the upward image of the
prescribed $E_i$-class; if the
corresponding sorts are NAMED
this upward map is a named
original unary coordinate map.

If instead $H\nleq J$,
then $H\cap J<H$.
The FULL NR2 hypothesis says the
relative index $[H:H\cap J]$
is INFINITE, giving infinitely
many distinct $J$-classes in $X$.
For simultaneous avoidance let
$G_X\le\Aut(\M)$ be the SETWISE
stabiliser of the actual $H$-class $X$.
It acts transitively on $X$:
for any $x,y\in X$, original
sort transitivity sends $x$ to
$y$, and by invariance of $H$
that automorphism preserves $X$.
For each unforced $J_k$, the
orbit of $[x]_{J_k}$ under $G_X$
has size exactly
$[H:H\cap J_k]=\infty$.
Apply \cref{lem:finite-action-avoidance} to the $G_X$-action on
$\bigsqcup_{k=1}^s M_a/J_k$, keeping the index $k$ as a sort tag.
The moving finite set consists of the tagged classes $(k,[x]_{J_k})$;
the forbidden finite set is the union of the tagged forbidden lists.
The preceding index calculation makes each moving orbit infinite.
Thus one $g\in G_X$ avoids all lists simultaneously, and $gx\in X$
is the required common real witness. This also covers $s=0$.
No full fine SPOKE colours to
additional real anchors are
controlled by this argument.
\end{proof}

\begin{corollary}[One fresh real frame for a compatible family of MI orphans]
\label{cor:typed-compatible-mi-single-real-frame}
In the setting of
\cref{lem:typed-mi-fibre-infinite-dichotomy},
suppose the $E_i$ are proper represented
MEET-IRREDUCIBLES on the SAME original
real sort $a$, with pairwise compatible
actual quotient class vertices
$b_i=C_i\in M_a/E_i$.
Put $A=\cl(\{b_1,\ldots,b_m\})$
in the ACTUAL MI-only hidden language.
For EVERY finite induced unary-closed
hidden $B\supseteq A$, there is ONE
original real point $x$ with
\[
   x\in\bigcap_i C_i,\qquad
   A\subseteq\cl(\{x\}),\qquad
   \cl(\{x\})\cap B=A.
 \tag{CMF}\label{eq:typed-compatible-mi-single-real-frame}
\]
In particular $x$ can avoid every
original real vertex of $B$.
This does NOT replace the full
\cref{lem:typed-fresh-representatives},
which allows mutually INCOMPATIBLE
or different-original-sort orphan bases.
\end{corollary}
\begin{proof}
The joint real fibre $X=\bigcap_i C_i$
is an infinite actual $H$-class,
$H=\bigcap_i E_i$, by
\cref{lem:typed-mi-fibre-infinite-dichotomy}.
Every $x\in X$ has each actual
$b_i$ as a NAMED quotient term
value, so
$A\subseteq\cl(\{x\})$.

There are only FINITELY many
real-source unary term functions
$\tau$ on sort $a$ by
\cref{lem:typed-unary-closure-bound}.
The kernel $K_\tau$ is either
actual real equality or a proper
represented MI by
\cref{lem:typed-coordinate-normal-form}.
If $H\le K_\tau$, then
$K_\tau\ne\Id$ because $H>\Id$.
The meet-prime implication
in \cref{lem:typed-mi-fibre-infinite-dichotomy}
provides some $E_i\le K_\tau$.
Thus the value $\tau(x)$ is
a unary upward-quotient/transition
term of $b_i$, CONSTANT on $X$,
and is already in the closed $A$.
This value is FORCED, not forbidden.

For every remaining term with
$H\nleq K_\tau$, the same lemma
gives an INFINITE set of possible
$\tau$-values on $X$, and
simultaneous finite-coset avoidance
lets us choose ONE $x\in X$
such that $\tau(x)$ avoids
EVERY hidden vertex of $B$
of its output sort, for ALL
such terms at once. In
particular the real-identity
term has kernel $\Id$, is
unforced and avoids $B_{\mathrm{real}}$.
The full hidden term closure
of $x$ consists of precisely
these finitely many forced and
unforced term values. Therefore
$\cl(\{x\})\cap B$ consists
exactly of the forced values
already in $A$, while the
reverse inclusion holds because
all $b_i$ are values of $x$.
\end{proof}

\begin{lemma}[Typed fresh representatives]
\label{lem:typed-fresh-representatives}
Let $A\le B$ be finite actual induced unary-closed hidden structures,
and let $F$ be any additional finite set of actual hidden vertices.
For every proper coordinate vertex $b\in A$ of sort $B_{a,J}$,
one can choose a real representative $x_b\in M_a$ with $F_J(x_b)=b$
such that all representatives are distinct and outside
$(B\cup F)_{\mathrm{real}}$. For every real-source term $\tau$ on $a$:
if $J\le K_\tau$, its forced value $\tau(x_b)$ is the prescribed
named term value of $b$ in $A$; otherwise $\tau(x_b)$ avoids
$B\cup F$ and every unforced term value of earlier representatives.
This gives literal closure freshness as well as fresh real representatives. With $C=A_{\mathrm{real}}$ and
$X=\{x_b:b\in A\text{ is a proper coordinate vertex}\}$, one has
\[
                  B\cap(C\cup X)^{\coord}=A.
\]
\end{lemma}
\begin{proof}
Work in one actual ambient copy of $\widehat\M$ and put $G=\Aut(\M)$.
All named projections and terms are $G$-equivariant. For the source sort
of $b$, \cref{lem:typed-unary-closure-bound,lem:typed-coordinate-normal-form}
give finitely many term functions, with kernel either real equality
(for the identity term) or an actual proper represented MI equivalence.
We verify the two inputs of
\cref{lem:foundational-quotient-fibre-avoidance}.

If $J\le K_\tau$, \cref{cor:typed-forced-term-factorisation}
gives a literal named factorisation $\tau=\sigma_{\tau,J}\circ F_J$.
The forced value $\sigma_{\tau,J}(b)$ belongs to $A$ by unary closure
and is retained. If $J\nleq K_\tau$, then $J\cap K_\tau<J$,
and FULL original NR2 makes $[J:J\cap K_\tau]$ infinite.
These are statements about the actual original invariant kernels,
including $K_\tau=\Id_a$, not merely about the named MI sublattice.
Original sort transitivity supplies the remaining group-action input.

Apply the common fibre-avoidance lemma successively to each $b$,
using as forbidden set $B\cup F$ together with all previously chosen
unforced term values. This is finite, and the lemma keeps every forced
value while avoiding the whole forbidden set on every unforced term.
Since $J>\Id_a$, the real identity term is unforced; it makes the
representatives distinct and avoids all old forbidden real points.
These are actual representatives of the same quotient vertices, so
all their prescribed upward maps and quotient transitions stay fixed.
The semantic-cylinder relations between retained hidden vertices
stay unchanged by \cref{lem:typed-cylinder-realisation}.

Every vertex of $A$ is either in $C$ or is a named projection of its
chosen $x_b$, giving $A\subseteq(C\cup X)^{\coord}$.
Conversely unary closure is the union of the one-generated closures.
Terms on $C$ stay in $A$; forced terms on each $x_b$ stay in $A$;
unforced terms avoid $B$. This proves the displayed intersection.
\end{proof}

\begin{remark}[CCO-free kernel criterion for fresh real frames]
\label{rem:typed-fresh-frame-cco-free}
The common provider \cref{lem:foundational-quotient-fibre-avoidance}
uses only a transitive group action, finitely many equivariant term
functions, literal named factorisation of forced terms and infinite
relative indices for unforced terms. Thus the finite freshness step
requires no original CCO, SAP, symmetry, lattice distributivity or
model homogeneity once these inputs have been verified.
In the actual terminal signature the factorisation is
\cref{cor:typed-forced-term-factorisation}, and the index condition is
FULL original NR2. CCO is used upstream to establish the represented
quotient geometry and real-base orbit infinitude. Fresh frames alone
do not imply hereditary hidden SAP: the real-base orbit condition is
still needed in \cref{lem:foundational-real-frame-orbits}.
Nor does this argument authorize adding non-MI meet quotients to
the hidden language without rechecking both the term-kernel conditions
and the real-base orbit hypothesis.
\end{remark}
